\documentclass[11pt]{article}
\usepackage{amsmath,amssymb,amsthm,mathtools,enumitem,geometry}
\geometry{margin=1in}
\newtheorem{theorem}{Theorem}
\newtheorem{lemma}{Lemma}
\newtheorem{proposition}{Proposition}
\newtheorem{definition}{Definition}
\newtheorem{remark}{Remark}
\newtheorem{corollary}{Corollary}
\DeclareMathOperator{\Ran}{Ran}
\DeclareMathOperator{\tr}{tr}
\DeclareMathOperator{\Fix}{Fix}
\newcommand{\Kcmp}{\mathsf K_{\rm cmp}}
\newcommand{\AZ}{\mathsf A_Z}
\newcommand{\Eef}{E_{\rm EF}}
\newcommand{\Htest}{\mathcal H_{\rm test}}
\newcommand{\Ecmp}{\mathcal E_{\rm cmp}}

\title{Step 71: The Weil--Douglas Bridge Gate for the RH Membrane Route}
\author{Working note}
\date{}

\begin{document}
\maketitle

\section{Purpose}
The formed-layer membrane framework reduces the RH-facing confinement route to
\[
  \AZ\preceq \Kcmp\preceq \Theta^- ,
  \qquad \tr\Theta^-\to0,
\]
with a separating anti-invariant zero readout and a fixed or exhaustive zero ledger.
The first load-bearing construction obligation is the bridge
\[
  \AZ\preceq\Kcmp.
\]
This note isolates the exact operator-theoretic content of that obligation.  It shows that the desired bridge is not a scalar explicit-formula identity and not a trace equality.  It is a Douglas contractive factorization gate.

\section{Zero-side and carrier-side maps}
Let $Y^-$ denote the anti-invariant response space for the completed duality.  The zero-side anti-invariant ledger is represented by a map
\[
   V_Z:Y^-\longrightarrow \mathcal Z,
\]
so that
\[
   \AZ=V_Z^*V_Z.
\]
The completed carrier-side feature map is required to have the form
\[
   W_{\rm cmp}:Y^-\longrightarrow \Ecmp,
   \qquad
   \Kcmp=W_{\rm cmp}^*W_{\rm cmp}.
\]
In the RH explicit-formula intended application one wants a visible decomposition
\[
   W_{\rm cmp}=\begin{bmatrix}
       W_{\rm pr}\\ W_\infty\\ W_{\rm pole}\\ W_{\rm tail}
   \end{bmatrix},
\]
with
\[
   \Kcmp=K_{\rm pr}+K_\infty+K_{\rm pole}+K_{\rm tail}.
\]
The terms are named prime-power, archimedean/gamma, pole/completion, and support/tail features.  This is a positive feature decomposition only after each block has been represented as a positive quadratic feature map.

\section{Douglas bridge theorem}
\begin{theorem}[Weil--Douglas bridge gate]
Let $V:Y\to\mathcal Z$ and $W:Y\to\mathcal E$ be bounded operators between Hilbert spaces.  The following are equivalent:
\begin{enumerate}[label=(\roman*)]
\item $V^*V\preceq W^*W$.
\item $\|Vy\|\le \|Wy\|$ for every $y\in Y$.
\item There exists a contraction $T:\overline{\Ran W}\to \overline{\Ran V}$ such that
\[
   V=TW.
\]
Equivalently, after extending $T$ by zero on $\overline{\Ran W}^{\perp}$, there is a contraction $T:\mathcal E\to\mathcal Z$ with $V=TW$.
\end{enumerate}
\end{theorem}

\begin{proof}
The equivalence of (i) and (ii) is immediate from
\[
   \|Vy\|^2=\langle y,V^*Vy\rangle,
   \qquad
   \|Wy\|^2=\langle y,W^*Wy\rangle.
\]
Assume (ii).  Define $T_0$ on $\Ran W$ by $T_0(Wy)=Vy$.  If $Wy=Wy'$, then $W(y-y')=0$, so (ii) gives $V(y-y')=0$; hence $T_0$ is well-defined.  Moreover $\|T_0Wy\|=\|Vy\|\le\|Wy\|$, so $T_0$ is a contraction on $\Ran W$ and extends uniquely to $\overline{\Ran W}$.  This gives (iii).  Conversely, if $V=TW$ with $\|T\|\le1$, then $\|Vy\|\le\|Wy\|$, hence (ii) and (i).
\end{proof}

\begin{corollary}[RH bridge target]
The explicit-formula domination obligation
\[
   \AZ\preceq\Kcmp
\]
is equivalent to constructing a contraction
\[
   T:\Ecmp\to\mathcal Z,
   \qquad
   V_Z=T W_{\rm cmp}.
\]
Thus the bridge is a contractive factorization, not merely equality of scalar trace terms.
\end{corollary}

\section{From a completed positive form to features}
\begin{definition}[Completed anti-invariant carrier form]
A completed anti-invariant carrier form is a positive semidefinite sesquilinear form
\[
   Q_{\rm cmp}:Y^-\times Y^-\to\mathbb C
\]
with a visible decomposition
\[
   Q_{\rm cmp}=Q_{\rm pr}+Q_\infty+Q_{\rm pole}+Q_{\rm tail}
\]
whenever the four terms are individually positive or have been grouped into positive feature blocks.  The associated currency is the operator $\Kcmp$ satisfying
\[
   Q_{\rm cmp}(y,y')=\langle y,\Kcmp y'\rangle.
\]
\end{definition}

\begin{proposition}[Kolmogorov feature realization]
If $Q_{\rm cmp}\succeq0$ on $Y^-$, then there exists a Hilbert space $\Ecmp$ and a map $W_{\rm cmp}:Y^-\to\Ecmp$ such that
\[
   Q_{\rm cmp}(y,y')=\langle W_{\rm cmp}y,W_{\rm cmp}y'\rangle.
\]
If $Q_{\rm cmp}=\sum_\alpha Q_\alpha$ with each $Q_\alpha\succeq0$, then one may take
\[
   W_{\rm cmp}=\bigoplus_\alpha W_\alpha.
\]
\end{proposition}

\begin{proof}
Take the quotient of $Y^-$ by the null space of $Q_{\rm cmp}$ and complete it in the induced inner product.  The map $W_{\rm cmp}$ sends $y$ to its equivalence class.  For a positive sum, apply the construction to each $Q_\alpha$ and concatenate the feature maps.
\end{proof}

\begin{theorem}[Positive Weil form implies Douglas gate]
Suppose a completed carrier form $Q_{\rm cmp}\succeq0$ satisfies
\[
   \|V_Z y\|^2\le Q_{\rm cmp}(y,y)
   \qquad\forall y\in Y^-.
\]
Let $W_{\rm cmp}$ be a feature realization of $Q_{\rm cmp}$.  Then there exists a contraction $T$ such that
\[
   V_Z=T W_{\rm cmp}.
\]
Equivalently,
\[
   \AZ\preceq\Kcmp.
\]
\end{theorem}

\begin{proof}
The inequality is exactly $V_Z^*V_Z\preceq W_{\rm cmp}^*W_{\rm cmp}$.  The result follows from the Weil--Douglas bridge gate.
\end{proof}

\section{Signed explicit formula terms and the completion defect}
The classical explicit formula is naturally signed: prime-power, gamma, pole, and tail terms enter with different signs depending on the convention and test class.  The membrane route does not accept a signed accounting identity as a positive feature map.

\begin{definition}[Completion defect]
Given a visible partial feature map $W_{\rm vis}$ and zero readout $V_Z$, define the explicit-formula defect $\Eef\succeq0$ by an admissible inequality
\[
   V_Z^*V_Z\preceq W_{\rm vis}^*W_{\rm vis}+\Eef.
\]
If $\Eef=0$, the bridge is exact.  If $\Eef\ne0$, it is a quantified allowance for missing gamma, pole, endpoint, support, tail, or positivity-completion data.
\end{definition}

\begin{proposition}[Defected Douglas bridge]
If
\[
   V_Z=T W_{\rm vis}+R,
   \qquad \|T\|\le1,
   \qquad R^*R\preceq\Eef,
\]
then for every $s>0$,
\[
   \AZ\preceq (1+s)W_{\rm vis}^*W_{\rm vis}+(1+s^{-1})\Eef.
\]
\end{proposition}

\begin{proof}
Use $(a+b)^*(a+b)\preceq (1+s)a^*a+(1+s^{-1})b^*b$ with $a=TW_{\rm vis}$ and $b=R$, then $W_{\rm vis}^*T^*TW_{\rm vis}\preceq W_{\rm vis}^*W_{\rm vis}$.
\end{proof}

\section{Why trace equality is not enough}
\begin{proposition}[Trace shadow counterexample]
There exist $A,K\succeq0$ with
\[
   \tr A=\tr K
\]
but
\[
   A\npreceq K.
\]
For example,
\[
   A=\begin{pmatrix}2&0\\0&1/2\end{pmatrix},
   \qquad
   K=\begin{pmatrix}1&0\\0&3/2\end{pmatrix}
\]
have equal trace, but $K-A$ has a negative eigenvalue.
\end{proposition}

\begin{remark}
This is the finite-dimensional reason the explicit formula cannot be used merely as total prime/zero trace balance.  The RH membrane route needs Loewner domination in the anti-invariant sector, equivalently a contractive factorization.
\end{remark}

\section{The first construction target}
The actual RH construction target can now be stated without framework ambiguity.

\begin{quote}
Construct an anti-invariant completed test/response space $Y^-$, a zero-side map $V_Z$, and a completed feature map
\[
   W_{\rm cmp}=W_{\rm pr}\oplus W_\infty\oplus W_{\rm pole}\oplus W_{\rm tail}
\]
from the Weil/Selberg explicit formula such that
\[
   V_Z=T W_{\rm cmp},\qquad \|T\|\le1.
\]
\end{quote}

This is stronger than writing the explicit formula.  It requires the signed formula to be reorganized as a positive operator-valued feature representation on the anti-invariant response family.

\section{Nonclaim boundary}
This note does not prove RH.  It does not construct the zeta carrier.  It does not prove Weil positivity in a new form.  It isolates the exact operator gate that a Weil-style construction must satisfy: Douglas-contractive domination of zero-side anti-invariant displacement by completed carrier features.

\end{document}
